% ============================================================================
%  Problems: ทฤษฎีจำนวน (Number Theory)
%  Ordered from easy (basic) → intermediate.
%  Shared by exercise.tex and answer-key.tex
% ============================================================================

% --- Part 1: พื้นฐาน (Basic) ---
\examsection{ตอนที่ 1: พื้นฐาน — การหารลงตัวและจำนวนเฉพาะ}

\problem{
  จงตรวจสอบว่า $72$ หารด้วย $6$ ลงตัวหรือไม่ พร้อมให้เหตุผล
}{
  $72 \div 6 = 12$ เป็นจำนวนเต็ม \quad ดังนั้น $6 \mid 72$

  หรือใช้กฎ: $72$ หารด้วย $2$ ลงตัว (หลักหน่วย $2$) และหารด้วย $3$ ลงตัว ($7+2=9$)
  ดังนั้นหารด้วย $6$ ลงตัว
}

\problem{
  จงแยกตัวประกอบเฉพาะของ $120$
}{
  $120 = 2 \times 60 = 2 \times 2 \times 30 = 2 \times 2 \times 2 \times 15$

  $= 2^3 \times 3 \times 5$
}

\problem{
  จงตรวจสอบว่า $127$ เป็นจำนวนเฉพาะหรือไม่
}{
  $\sqrt{127} \approx 11.3$ — ทดลองหารด้วยจำนวนเฉพาะ $\leq 11$:

  $2 \nmid 127$, $3 \nmid 127$ ($1+2+7=10$), $5 \nmid 127$, $7 \nmid 127$,
  $11 \nmid 127$

  ไม่มีตัวใดหารลงตัว $\rightarrow$ $127$ เป็นจำนวนเฉพาะ
}

% --- Part 2: การประยุกต์ (Intermediate) ---
\examsection{ตอนที่ 2: การประยุกต์ — GCD/LCM และสมภาค}

\problem{
  จงหา ห.ร.ม. และ ค.ร.น. ของ $48$ และ $180$ โดยวิธีการแยกตัวประกอบ
}{
  $48 = 2^4 \times 3$ \qquad $180 = 2^2 \times 3^2 \times 5$

  GCD $= 2^{\min(4,2)} \times 3^{\min(1,2)} = 2^2 \times 3 = 12$

  LCM $= 2^{\max(4,2)} \times 3^{\max(1,2)} \times 5 = 2^4 \times 3^2 \times 5 = 720$
}

\problem{
  จงหา ห.ร.ม. ของ $144$ และ $89$ โดยใช้ขั้นตอนวิธีของยุคลิด
}{
  $144 = 1 \times 89 + 55$ \quad $\rightarrow$ \quad GCD(89, 55)

  $89 = 1 \times 55 + 34$ \quad $\rightarrow$ \quad GCD(55, 34)

  $55 = 1 \times 34 + 21$ \quad $\rightarrow$ \quad GCD(34, 21)

  $34 = 1 \times 21 + 13$ \quad $\rightarrow$ \quad GCD(21, 13)

  $21 = 1 \times 13 + 8$ \quad $\rightarrow$ \quad GCD(13, 8)

  $13 = 1 \times 8 + 5$ \quad $\rightarrow$ \quad GCD(8, 5)

  $8 = 1 \times 5 + 3$ \quad $\rightarrow$ \quad GCD(5, 3)

  $5 = 1 \times 3 + 2$ \quad $\rightarrow$ \quad GCD(3, 2)

  $3 = 1 \times 2 + 1$ \quad $\rightarrow$ \quad GCD(2, 1)

  $2 = 2 \times 1 + 0$ \quad $\rightarrow$ \quad GCD = 1
}

\problem{
  จงหาค่าของ $17 \times 23 \pmod{7}$ โดยใช้สมบัติของสมภาค
}{
  $17 \equiv 3 \pmod{7}$ \quad ($17 = 2 \times 7 + 3$)

  $23 \equiv 2 \pmod{7}$ \quad ($23 = 3 \times 7 + 2$)

  $17 \times 23 \equiv 3 \times 2 = 6 \pmod{7}$
}

\problem{
  จงหาหลักหน่วยของ $3^{2024}$ (หรือเศษจากการหาร $3^{2024}$ ด้วย $10$)
}{
  $3^n \bmod 10$: วัฏจักร $\{3, 9, 7, 1\}$ \quad ($3^1=3$, $3^2=9$, $3^3=27\equiv 7$, $3^4=81\equiv 1$)

  $2024 \div 4 = 506$ เศษ $0$ \quad $\rightarrow$ \quad $3^{2024} \equiv 3^4 \equiv 1 \pmod{10}$

  หลักหน่วยคือ $1$
}

\problem{
  จงหาเศษจากการหาร $2^{100}$ ด้วย $5$
}{
  ใช้ Fermat: $2^4 \equiv 1 \pmod{5}$ \quad ($p=5$ เป็นจำนวนเฉพาะ)

  $100 = 4 \times 25$ \quad $\rightarrow$ \quad
  $2^{100} = (2^4)^{25} \equiv 1^{25} \equiv 1 \pmod{5}$
}

\problem{
  วันที่ 1 มกราคม 2025 เป็นวันพุธ อยากทราบว่าวันที่ 31 ธันวาคม 2025
  (อีก $364$ วัน) จะเป็นวันอะไร?
}{
  $364 \div 7 = 52$ เศษ $0$

  เศษ $0$ หมายถึงวนกลับมาที่เดิม

  ดังนั้นวันที่ 31 ธันวาคม 2025 เป็น วันพุธ เช่นกัน
}

\problem{
  จงหาหลักหน่วยของ $7^{2026} + 9^{2026}$
}{
  $7^n \bmod 10$: วัฏจักร $\{7, 9, 3, 1\}$ \quad
  $2026 \div 4 = 506$ เศษ $2$ \quad $\rightarrow$ \quad $7^{2026} \equiv 9$

  $9^n \bmod 10$: วัฏจักร $\{9, 1\}$ (สลับ $9$ กับ $1$) \quad
  $2026$ เป็นเลขคู่ \quad $\rightarrow$ \quad $9^{2026} \equiv 1$

  $9 + 1 = 10 \equiv 0 \pmod{10}$

  หลักหน่วยคือ $0$
}

% --- Part 3: ระดับ สอวน. (POSN Level) ---
\examsection{ตอนที่ 3: ระดับข้อสอบ สอวน.}

\problem{
  (ข้อสอบจริงปี 68 ข้อ 3) ตัวเลขหลักหน่วย ของผลคูณนี้ $(2024)^{123} \cdot (2025)^{234} \cdot (2026)^{345}$ คือเลขใด
}{
  $0$
}

\problem{
  (ข้อสอบจริงปี 68 ข้อ 4) สำหรับจำนวนเต็ม $a$ ใดๆ ตัวเลขหลักหน่วยของ $a^4$ ที่เป็นไปได้ทั้งหมดมีกี่ตัว  
}{
  $4$
}

\problem{
  (ข้อสอบจริงปี 68 ข้อ 5) จงหาผลรวมของจำนวนเต็มบวก $x$ ทุกจำนวนที่ $x<20$ \\
  และสอดคล้องกับ $2x+3 \equiv 4\pmod{11}$
}{
  $23$
}

\problem{
  (ข้อสอบจริงปี 67 ข้อ 12) สำหรับจำนวนเต็มบวก $a,x$ และ $y$ ข้อใดกล่าว \textbf{ไม่ถูกต้อง} 
  \begin{enumerate}
    \item[ก.] ถ้า $x$ และ $y$ ต่างก็หารด้วย $a$ ไม่ลงตัว แล้ว $(x+y)$ จะหารด้วย $a$ ไม่ลงตัว 
    \item[ข.] ถ้า $x$ และ $y$ ต่างก็หารด้วย $a$ ลงตัว แล้ว $(x-y)$ จะหารด้วย $a$ ลงตัว 
    \item[ค.] ถ้า $x$ และ $y$ ต่างก็หารด้วย $a$ ลงตัว แล้ว $(xy)$ จะหารด้วย $a$ ลงตัว 
    \item[ง.] ถ้า $x$ หารด้วย $a$ ลงตัว และ $y$ หารด้วย $a$ ไม่ลงตัว แล้ว $(x+y)$ จะหารด้วย $a$ ไม่ลงตัว 
  \end{enumerate}
}{
  ก.
}

\problem{
  (ข้อสอบจริงปี 67 ข้อ 37) โรงเรียนเสริมพัฒนาการแห่งหนึ่งมี 2 สาขา ต้องการจัดกิจกรรมให้กับนักเรียนในแต่ละสาขาในเวลาที่ต่างกัน โดยโรงเรียนสาขา A มีนักเรียน 6,786 คน ส่วนโรงเรียนสาขา B มีนักเรียน 13,260 คน ทางโรงเรียนเสริมพัฒนาการ ต้องการแบ่งกลุ่มให้นักเรียนในแต่ละสาขา มีจำนวนสมาชิกเป็นเลขคี่และเท่ากันทุกกลุ่ม จงหาจำนวนคี่ที่มากที่สุดที่จะแบ่งกลุ่มนักเรียนได้พอดี 
}{
  $39$
}

\problem{
  (ข้อสอบจริงปี 66 ข้อ 43) จงหาเศษที่เหลือจากการหาร $1^{2566}+2^{2566}+3^{2566}+4^{2566}$ ด้วย $10$
}{
  $0$
}

\problem{
  (ข้อสอบจริงปี 65 ข้อ 43) ค.ร.น. ของ $2^33^57^2$ และ $2^43^3$ ตรงกับข้อใด (ตอบในรูปผลคูณของตัวประกอบเฉพาะยกกำลัง)
}{
  $2^43^57^2$
}